\section{Tohumlar Arası Değişkenlik}
\label{app:std}

Tablo~\ref{tab:std}, Tablo~\ref{tab:main}'deki her öğrenen sistem için üç tohum üzerinden ortalama $\pm$ örneklem standart sapmasını verir. mDeBERTa varyantlarında standart sapma en fazla 2.9 puandır (medyan yaklaşık 1.2); bu, ortak çözümlemenin ortalama HardCons kazancının ($+3.0$) altındadır. XLM-R çok daha kararsızdır (20 puana kadar); bunun başlıca nedeni çöken İspanyolca tohumudur.

\begin{table*}[t]
\centering\small
\setlength{\tabcolsep}{3pt}
\resizebox{\textwidth}{!}{%
\begin{tabular}{llcccccccccccc}
\toprule
 & & \multicolumn{3}{c}{EN} & \multicolumn{3}{c}{ES} & \multicolumn{3}{c}{EU} & \multicolumn{3}{c}{Karma} \\
Model & Çöz. & F1 & Soft & Hard & F1 & Soft & Hard & F1 & Soft & Hard & F1 & Soft & Hard \\
\midrule
XLM-R & ayrı & 69.1$\pm$1.6 & 82.7$\pm$1.3 & 72.1$\pm$2.2 & 54.2$\pm$10.7 & 66.1$\pm$15.2 & 51.9$\pm$16.1 & 61.1$\pm$3.8 & 80.5$\pm$2.0 & 63.7$\pm$0.7 & 63.4$\pm$0.2 & 79.3$\pm$1.2 & 64.9$\pm$0.7 \\
 & ortak & 69.6$\pm$1.6 & 93.9$\pm$0.3 & 77.3$\pm$1.1 & 58.6$\pm$5.7 & 93.5$\pm$4.1 & 68.4$\pm$2.8 & 61.0$\pm$3.0 & 94.3$\pm$2.7 & 70.0$\pm$1.2 & 63.8$\pm$0.1 & 88.4$\pm$1.3 & 69.7$\pm$0.7 \\
mDeBERTa & ayrı & 79.6$\pm$0.9 & 87.1$\pm$1.7 & 81.0$\pm$0.6 & 77.3$\pm$0.6 & 85.6$\pm$0.8 & 78.3$\pm$0.3 & 76.5$\pm$0.7 & 82.7$\pm$2.3 & 75.7$\pm$1.9 & 73.4$\pm$1.0 & 81.8$\pm$1.6 & 73.8$\pm$1.5 \\
 & ortak & 80.3$\pm$1.3 & 91.8$\pm$0.6 & 84.5$\pm$1.0 & 77.9$\pm$0.7 & 89.3$\pm$0.9 & 81.1$\pm$1.2 & 76.7$\pm$0.4 & 87.5$\pm$2.1 & 78.8$\pm$1.2 & 73.5$\pm$0.9 & 86.2$\pm$1.0 & 76.6$\pm$1.0 \\
mDeBERTa-NLI & ayrı & 82.8$\pm$0.7 & 88.9$\pm$0.6 & 84.6$\pm$0.9 & 79.0$\pm$0.8 & 85.0$\pm$0.7 & 79.5$\pm$0.7 & 75.1$\pm$1.7 & 81.8$\pm$2.3 & 75.1$\pm$2.1 & 74.7$\pm$0.9 & 82.4$\pm$1.8 & 74.7$\pm$1.6 \\
 & ortak & 83.0$\pm$0.9 & 94.2$\pm$0.3 & 87.5$\pm$1.2 & 78.7$\pm$0.7 & 90.4$\pm$0.5 & 82.2$\pm$0.6 & 75.6$\pm$1.3 & 86.8$\pm$1.1 & 78.5$\pm$1.5 & 74.8$\pm$0.8 & 88.4$\pm$1.5 & 78.3$\pm$1.4 \\
mDeBERTa+$\mathcal{L}_\text{cons}$ & ayrı & 80.7$\pm$0.4 & 90.0$\pm$2.3 & 83.9$\pm$1.6 & 76.5$\pm$1.0 & 86.9$\pm$2.0 & 78.6$\pm$1.5 & 75.1$\pm$0.3 & 85.4$\pm$1.2 & 77.4$\pm$1.0 & 72.4$\pm$0.5 & 85.7$\pm$1.2 & 75.4$\pm$0.9 \\
 & ortak & 80.8$\pm$0.4 & 92.5$\pm$0.7 & 85.3$\pm$1.2 & 76.8$\pm$1.5 & 91.1$\pm$1.6 & 81.2$\pm$1.8 & 75.3$\pm$0.3 & 88.0$\pm$0.7 & 79.3$\pm$0.7 & 72.5$\pm$0.4 & 88.6$\pm$1.3 & 77.2$\pm$0.9 \\
\bottomrule
\end{tabular}}
\caption{Dev sonuçları (\%), üç tohum üzerinden ortalama $\pm$ standart sapma.}
\label{tab:std}
\end{table*}


\section{Diller Arası Aktarımın Tüm Sonuçları}
\label{app:transfer}

Tablo~\ref{tab:transfer_full}, Tablo~\ref{tab:transfer}'ü SoftCons, ayrı çözümleme ve standart sapmalarla genişletir. Ortak çözümleme HardCons'u her eğitim ayarında ve her track'te artırır; en büyük artış sıfır-örnek ayarındadır.

\begin{table*}[t]
\centering\small
\setlength{\tabcolsep}{3pt}
\resizebox{\textwidth}{!}{%
\begin{tabular}{llcccccccccccc}
\toprule
 & & \multicolumn{3}{c}{EN} & \multicolumn{3}{c}{ES} & \multicolumn{3}{c}{EU} & \multicolumn{3}{c}{Karma} \\
Eğitim verisi & Çöz. & F1 & Soft & Hard & F1 & Soft & Hard & F1 & Soft & Hard & F1 & Soft & Hard \\
\midrule
Kendi verisi & ayrı & 79.6$\pm$0.9 & 87.1$\pm$1.7 & 81.0$\pm$0.6 & 77.3$\pm$0.6 & 85.6$\pm$0.8 & 78.3$\pm$0.3 & 76.5$\pm$0.7 & 82.7$\pm$2.3 & 75.7$\pm$1.9 & 73.4$\pm$1.0 & 81.8$\pm$1.6 & 73.8$\pm$1.5 \\
 & ortak & 80.3$\pm$1.3 & 91.8$\pm$0.6 & 84.5$\pm$1.0 & 77.9$\pm$0.7 & 89.3$\pm$0.9 & 81.1$\pm$1.2 & 76.7$\pm$0.4 & 87.5$\pm$2.1 & 78.8$\pm$1.2 & 73.5$\pm$0.9 & 86.2$\pm$1.0 & 76.6$\pm$1.0 \\
Yalnız EN & ayrı & 79.7$\pm$1.1 & 87.7$\pm$1.4 & 81.7$\pm$0.7 & 73.2$\pm$0.6 & 85.7$\pm$0.7 & 75.7$\pm$0.9 & 64.7$\pm$1.3 & 82.7$\pm$2.6 & 64.9$\pm$2.1 & 52.9$\pm$3.6 & 70.8$\pm$3.6 & 49.3$\pm$2.7 \\
 & ortak & 79.6$\pm$0.7 & 91.3$\pm$0.8 & 84.0$\pm$0.2 & 74.4$\pm$0.7 & 92.7$\pm$0.7 & 80.7$\pm$0.5 & 64.7$\pm$1.3 & 89.4$\pm$2.0 & 68.5$\pm$2.0 & 53.1$\pm$3.6 & 80.5$\pm$4.1 & 54.4$\pm$3.2 \\
Birleşik & ayrı & 80.4$\pm$0.6 & 88.9$\pm$1.0 & 82.8$\pm$1.0 & 78.2$\pm$1.1 & 87.0$\pm$0.8 & 80.1$\pm$0.9 & 76.6$\pm$0.6 & 85.6$\pm$1.1 & 78.9$\pm$0.5 & 73.9$\pm$0.6 & 83.0$\pm$0.3 & 74.5$\pm$0.7 \\
 & ortak & 80.8$\pm$0.5 & 91.6$\pm$0.6 & 84.8$\pm$0.9 & 78.7$\pm$1.0 & 91.8$\pm$1.8 & 83.6$\pm$1.2 & 77.0$\pm$1.0 & 89.3$\pm$1.3 & 81.5$\pm$1.2 & 74.2$\pm$0.5 & 87.8$\pm$0.2 & 77.7$\pm$0.7 \\
\bottomrule
\end{tabular}}
\caption{mDeBERTa'nın diller arası aktarımı (dev, \%), üç tohum üzerinden ortalama $\pm$ std., ayrı ve ortak çözümlemeyle. Eğitim verileri Tablo~\ref{tab:transfer}'teki gibidir.}
\label{tab:transfer_full}
\end{table*}


\section{Hiperparametreler ve Uygulama}
\label{app:hparams}

\begin{itemize}\itemsep0pt
\item Modeller (Hugging Face): \texttt{FacebookAI/xlm-roberta-base}, \texttt{microsoft/mdeberta-v3-base} ve \texttt{MoritzLaurer/mDeBERTa-v3-base-xnli-}\allowbreak\texttt{multilingual-nli-2mil7}.
\item Öğrenme oranı (EN doğrulama diliminde $\{2,5\}\cdot10^{-5}$ arasından seçildi): XLM-R $2\cdot10^{-5}$, mDeBERTa $5\cdot10^{-5}$, mDeBERTa-NLI $5\cdot10^{-5}$.
\item Epoch: en fazla 10 (tek dilli, yalnız EN) ya da 4 (karma, birleşik); en iyi epoch, ayrı çözümlemeyle doğrulama ağırlıklı F1'ine göre seçilir.
\item Optimizasyon: AdamW, 32'lik yığın (bir çiftin iki yönü her zaman aynı yığında), \%10 doğrusal ısınma ve ardından doğrusal azalma, 0.01 ağırlık azaltımı, 1.0 gradyan kırpma, en fazla 64 token. Kelime gömmeleri dondurulmuştur (250 bin tokenlik sözlükler).
\item Tutarlılık kaybı ağırlığı $\lambda=1$ (ayarlanmadı).
\item Tohumlar 1, 2, 3; tohum toplulukları tahmin edilen olasılıkların ortalamasını alır. Apple M3 Pro'da tek dilli koşu başına 3.5--6 dakika.
\item Doğrulama dilimi: \texttt{pair\_id}'nin MD5 özeti en düşük \%10'a düşen kaynak çiftler (1{,}760 eğitim kaynak çiftinden 175'i).
\end{itemize}
